\sscp{Boano}{10-03-01}
Boano is spoken on a small island of the same name in the \ili{Manipa}
Straits off the northwest coast of Seram.
Many Boano speakers were forcibly relocated by the Dutch to \ili{Manipa}
Island in the mid 1600s at the end of the Hoamoal wars,
in efforts to gain a monopoly on the spice trade.

\citet{Stresemann1927} included \ili{Manipa}, \ili{Kelang},
and Boano within his Sub-\ili{Ambon} branch of Proto-\ili{Ambon}.
\citeauthor{Stresemann1927}’s approach to the apical consonants for his Sub-\ili{Ambon}
branch was to posit a merger of what we now call PMP *d/*z/*R/*l/*j
> **l, and then posit a complex set of patterned splits \citep[23--29]{Stresemann1927},
as discussed in \srf{10-01} above.

\citet{Collins1983} includes Boano within the \tsc{\ili{Piru} Bay}
branch of his “Nunusaku”,
with its nearest relatives being \ili{Wakasihu}, \ili{Alang}, and \ili{Larike}.
He sees the Boano data (\citeyear[88--93]{Collins1983})
as requiring a separate reflex
for *d/*z “which contrasts with the reflex for *R,
*l and *j.” Note that he sees a single reflex for the latter
three proto-phonemes, implying a merger.
\citet[88]{Collins1983} concedes that while Stresemann’s assumption
of a merger was incorrect,
nevertheless “At a later stage mergers did take place in many
languages and these changes are succinctly summarised
in Stresemann’s ‘Liquida Gesetze’ [Liquid laws].”
The “liquid laws” are summarised
and repackaged in \trf{10-01}.

We agree that Boano is an important witness for understanding
the phonological prehistory of our \tsc{\ili{Ambon}-Seram} subgroup.
Our interpretation of the Boano data
and sound correspondences diverges from both \citeauthor{Stresemann1927}
and \citeauthor{Collins1983}, but is a bit closer to that of \citeauthor{Collins1983}.

Boano belongs with \tsc{\ili{Ambon}-Seram} on the basis of the merger
of *n/*ŋ/*ñ > \it{n}; *mb > \it{p}; *p > {\0}; *k > {\0}
and  *kahiw > **kai.
We also note that *b > \it{h},
*w > {\0}, and *ai/*a(C)i > \it{e} (or \it{ai}),
and *au/*a(C)u > \it{o}, \it{ou}.

Based on the treatment of the diphthongs,
\citet[69--73]{Collins1983} groups Boano with \ili{Larike}
and \ili{Wakasihu}, two languages now relocated to the western end of \ili{Ambon} Island.
Example \qf{10-25} shows *ai/*a(C)i > \it{e},
with the exception of \it{tai-a} ‘faeces’ which could be influenced
from the local \ili{Malay} \it{tai},
or any number of nearby languages.
\\

\noindent{\wg\st{6pt}\begin{tabularx}
{\linewidth}{>{\hspace{-\tabcolsep}}l<{\hspace{-\tabcolsep}}
>{\hspace{-\tabcolsep}\enspace}lll
>{\raggedright\arraybackslash}X}
\lsptoprule
%\tbx{10-25}	&	PMP/PCM	&	Boano	&	gloss	&		\\	\midrule
%\endfirsthead
\tbx{10-25}	&	PMP/PCM	&	Boano	&	gloss	&	note	\\	\midrule
%\endhead
	&	*kahiw	&	\tbr{e}-niʔi	&	tree, wood 	&	Proto-AS **kai > **ke, cf. \ili{Manipa} \it{hen}	\\
	&	*wahiyəR	&	en-eʔe	&	water 	&	**wain > **wen > \it{en} [*R > \it{n}]; cf. \ili{Manipa} \it{wel-e}	\\
	&	*qaqay, *waqay	&	\tbr{ee}	&	foot, leg 	&	cf. \ili{Manipa} \it{ai}, \ili{Wakasihu} \it{wei}	\\
	&	**saqi	&	s\tbr{e}-ʔe	&	paddle 	&	cf. \ili{Manipa} \it{sai}, \ili{Wakasihu} \it{sei}	\\
	&	*taqi	&	t\tbr{ai}-a	&	faeces	&	\hp{\,}	\\
\lspbottomrule
\end{tabularx}\mbox{}\\}

%Example \qf{10-26} shows *au/*a(C)u > \it{o,
The examples in \trf{10-26} shows *au/*a(C)u > \it{o,
ou} with the exception of *hapuy > \it{au} ‘fire’ which could
indicate that *p > {\0} occurred later than *h > {\0},
and after *au/*a(C)u > \it{o}.

\begin{table}\tcp{Boano and related languages *a(C)u}{10-26}
%\begin{adjustbox}{width=\textwidth}
	\begin{threeparttable} %\st{2pt}
	\begin{tabular}{lllllll}\lsptoprule
local gl.	&	far	&	odour	&	mango	&	fire	&	full	&	brother-in-law	\\
PMP	&	*zauq	&	*bahu	&	*pahuq	&	*hapuy	&	**taqun\su{†}	&	**ndau\su{‡}	\\ \midrule
Boano	&	-l\tbr{ou}	&	{\hr}h\tbr{o}-ni	&	{\hr}\tbr{o}-eʔe	&	{\hr}\tbr{au}	&	{\hr\hr}t\tbr{o}n-iʔi	&	{\hr\hr}r\tbr{o}-ʔu	\\
\ili{Manipa}	&	-lau	&	{\hr}hon	&	{\hr}au	&	{\hr}au	&	{\hr\hr}u-kon-e	&	{\hr\hr}ro-	\\
\ili{Wakasihu}	&	{\hr}lou	&	{\hr}hou	&	{\hr}owo	&	{\hr}ao-o	&	{\hr\hr}toun-a	&	{\hr\hr}ro-ʔu	\\
\ili{Asilulu}	&	{\hr}lau	&	{\hr}hau-ni	&		&	{\hr}au	&		&		\\
\ili{Kelang}	&		&		&		&		&	{\hr\hr}taun-a	&	{\hr\hr}rau	\\
\ili{Luhu}	&		&		&	{\hr}au	&	{\hr}au	&	{\hr\hr}taun-a	&	{\hr\hr}rau	\\
		\lspbottomrule
	\end{tabular}
		\begin{tablenotes}\tnsize
			\item[†] {Possibly from *taqun `year'. cf. \ili{Buru} \it{taun} `full'. \ili{Manipa} has *t > \it{k}.}
			\item[‡] {cf. \ili{Buru} \it{dawe} `brother-in-law (m. s.)' }
		\end{tablenotes}
	\end{threeparttable}
%\end{adjustbox}
\end{table}

%\noindent{\wg\st{6pt}\begin{tabularx}
%{\linewidth}{>{\hspace{-\tabcolsep}}l<{\hspace{-\tabcolsep}}
%>{\hspace{-\tabcolsep}\enspace}lll
%>{\raggedright\arraybackslash}X}
%\lsptoprule
%%\tbx{10-26}	&	PMP/PCM	&	Boano	&	gloss	&	compare	\\	\midrule
%%\endfirsthead
%\tbx{10-26}	&	PMP/PCM	&	Boano	&	gloss	&	compare	\\	\midrule
%%\endhead
	%&	*zauq	&	-l\tbr{ou}	&	far	&	\ili{Manipa} \it{-lau}, \ili{Wakasihu} \it{lou}, \ili{Asilulu} \it{lau}	\\
	%&	*bahu	&	h\tbr{o}-ni	&	odour	&	\ili{Manipa} \it{hon}, \ili{Wakasihu} \it{hou}, \ili{Asilulu} \it{hau-ni}	\\
	%&	*pahuq	&	\tbr{o}-eʔe	&	mango	&	\ili{Manipa} \it{au}, \ili{Wakasihu} \it{owo}, \ili{Luhu} \it{au}	\\
	%&	**taun	&	t\tbr{o}n-iʔi	&	full	&	\ili{Manipa} \it{u-kon-e} [*t > k], \ili{Wakasihu} \it{toun-a}, \ili{Kelang}, \it{taun-a}, \ili{Luhu} \it{taun-a}, \ili{Asilulu} \it{taun}, \ili{Buru} \it{taun}	\\
	%&	**ndau	&	r\tbr{o}-ʔu	&	brother-in-law	&	\ili{Manipa} \it{ro-}, \ili{Wakasihu} \it{ro-ʔu}, \ili{Kelang} \it{rau}, \ili{Luhu} \it{rau}, \ili{Asilulu} \it{rau}, \ili{Buru} \it{dawe}	\\
	%&	*hapuy	&	\tbr{au}	&	fire	&	\ili{Manipa} \it{au}, \ili{Wakasihu} \it{ao-o}, \ili{Asilulu} \it{au}, \ili{Luhu} \it{au}	\\
%\lspbottomrule
%\end{tabularx}\mbox{}\\}

Turning to the complex correspondences of *d,
*z, *l, *R and *j, it is important to first note the (regressive
and progressive) assimilation frequently encountered in the Boano data.
There are several words with \it{n{\ldots}n}
and \it{l{\ldots} l}, but none in the available data that mix
\it{n{\ldots}l} or \it{l{\ldots}n}.
Assimilation of this sort is not unusual in the region.
For example, in \ili{Buru},
out of a corpus of over 7,000 words \citep{GrimesGrimes2020},
only one word mixes \it{r{\ldots}l},
a place name \it{Nam-role} ‘Role bay’.
Otherwise we find \it{r{\ldots}r} or \it{l{\ldots}l} through regressive assimilation.
Consistently in \ili{Buru}, PMP *d > \it{r} (over 80 instances),
but there are a few words such as *daləm > \it{lale-n} ‘inside;
seat of emotions, character and social relations’ which have *d > \it{l} through regressive assimilation.

Consistent with \tsc{\ili{Ambon}-Seram}, Boano merges PMP *d/*z > \it{l},
as seen in \qf{10-27} and \qf{10-28}.

\noindent{\wg\st{6pt}\tblr{>{\hspace{-\tabcolsep}}l<{\hspace{-\tabcolsep}}
>{\hspace{-\tabcolsep}\enspace}
lll}
\lsptoprule
%\tbx{10-27}	&	PMP	&	Boano	&	gloss	\\	\midrule
%\endfirsthead
	&	PMP	&	Boano	&	gloss	\\	\midrule
%\endhead
\tbx{10-27}	&	*\tbr{d}aya	&	\tbr{l}aa	&	land(ward), inland	\\
	&	*\tbr{d}atu	&	\tbr{l}atu	&	chief	\\
	&	*\tbr{d}əpa	&	\tbr{l}ea	&	fathom	\\
	&	*\tbr{d}uRi	&	\tbr{l}uli	&	fishbone, thorn	\\
	&	*\tbr{d}aRaq	&	\tbr{l}ala	&	blood	\\
	&	*\tbr{d}aləm	&	-\tbr{l}ale	&	in, inside	\\
	&	*\tbr{d}aki	&	\tbr{l}ai-a	&	dirt on skin, grime	\\
	&	*\tbr{d}uha	&	\tbr{l}ua	&	two	\\
	&	*\tbr{d}uyuŋ	&	\tbr{l}uin-e	&	dugong, sea cow	\\
	&	*ha\tbr{d}iRi	&	\tbr{l}ili	&	house post	\\
	&	*m-u\tbr{d}əhi	&	mu\tbr{l}i	&	behind	\\
	&	*kə\tbr{d}əŋ	&	e\tbr{l}e	&	stand, stretch	\\
	&	*ku\tbr{d}əŋ	&	u\tbr{l}en	&	clay cooking pot	\\
	&	*si\tbr{d}uk	&	si\tbr{l}u	&	PMP ‘ladle’, Boano ‘scoop net’	\\
	&	*tu\tbr{d}uk	&	-tu\tbr{l}u	&	PMP ‘leak, drip, as a leaky roof’, Boano ‘eaves’	\\	\midrule
\tbx{10-28}	&	*\tbr{z}alan	&	\tbr{l}alan-e	&	path	\\
	&	*\tbr{z}auq	&	-\tbr{l}ou	&	far	\\
	&	*qu\tbr{z}an	&	u\tbr{l}an-e	&	rain	\\
	&	*qa\tbr{z}ay	&	a\tbr{l}a-	&	jaw, chin	\\
	&	*haRə\tbr{z}an	&	e\tbr{l}an-e	&	stair (PCM **əzan)	\\
\lspbottomrule
\end{tabular}\mbox{}\\}

\newpage
There is a split of PMP *R > \it{n,
l}, {\0}, with the latter two being conditioned reflexes.
PMP *R > \it{n} initially
and medially, as shown in \qf{10-29}.
Final PMP *R{\#} > \it{n,} {\0}.
Progressive assimilation is shown in \qf{10-30},
in which *R > \it{l} /lV{\gap}.
The loss of final *R{\#} > {\0} in some items is shown in \qf{10-31}.

\noindent{\wg\st{6pt}\tblr{>{\hspace{-\tabcolsep}}l<{\hspace{-\tabcolsep}}
>{\hspace{-\tabcolsep}\enspace}
lll}
\lsptoprule
%\tbx{10-29}	&	PMP/PCM	&	Boano	&	gloss	\\	\midrule
%\endfirsthead
	&	PMP/PCM	&	Boano	&	gloss	\\	\midrule
%\endhead
\tbx{10-29}	&	*\tbr{R}usuk	&	\tbr{n}usu	&	rib, ribcage	\\
	&	*\tbr{R}umaq	&	\tbr{n}uma	&	house	\\
	&	*u\tbr{R}at	&	u\tbr{n}at-e	&	vein, tendon	\\
	&	**tubu\tbr{R}a	&	tuhu\tbr{n}a	&	spit, spray healing mixture	\\
	&	*ta\tbr{R}uq	&	ta\tbr{n}u-ʔe	&	put, place	\\
	&	*ku\tbr{R}ita	&	a-u\tbr{n}ita	&	octopus	\\
	&	*tabu\tbr{R}i	&	tahu\tbr{n}i	&	conch (trumpet)	\\
	&	*baqə\tbr{R}u	&	he\tbr{n}u	&	new	\\
	&	*ma-bu\tbr{R}uk	&	ma-hu\tbr{n}u	&	rotten	\\
	&	*tu\tbr{R}un	&	tu\tbr{n}u	&	descend	\\
	&	*ti\tbr{R}əm	&	ti\tbr{n}en-e	&	oyster	\\
	&	*ma\tbr{R}uqanay	&	me\tbr{n}eʔe	&	male	\\
	&	*qaba\tbr{R}a	&	ha\tbr{n}a	&	shoulder 	\\
	&	*wa\tbr{R}əj	&	a\tbr{n}e-te	&	vine, rope	\\
	&	*ma-bə\tbr{R}əqat	&	me-ha\tbr{n}a	&	heavy	\\
	&	*laya\tbr{R}	&	naa\tbr{n}	&	sail (n)	\\
	&	*wahiyə\tbr{R}	&	e\tbr{n}-eʔe	&	water (**wain > **wen) 	\\
	&	*niu\tbr{R} > **niwe\tbr{R} 	&	nie\tbr{n}-e	&	coconut	\\
	&	*busu\tbr{R}	&	husu\tbr{n}-e	&	bow	\\ \midrule
\tbx{10-30}	&	*du\tbr{R}i	&	lu\tbr{l}i	&	fishbone, thorn	\\
	&	*da\tbr{R}aq	&	la\tbr{l}a	&	blood	\\
	&	*hadi\tbr{R}i	&	li\tbr{l}i	&	house post	\\ \midrule
\tbx{10-31}	&	*dəŋə\tbr{R}	&	pama-nene	&	hear, listen	\\
	&	*butəli\tbr{R}	&	hutil-e	&	wart	\\
	&	*lasə\tbr{R}	&	lase	&	testicles	\\
\lspbottomrule
\end{tabular}\mbox{}\\}

\newpage
There remain four exceptions in the available Boano data in which *R > \it{l}.
There is also one possible instance of PMP *r > \it{l}.
These are shown in example \qf{10-32}.

\noindent{\wg\st{6pt}\tblr{>{\hspace{-\tabcolsep}}l<{\hspace{-\tabcolsep}}
>{\hspace{-\tabcolsep}\enspace}
lll}
\lsptoprule
%\tbx{10-32}	&	PMP	&	Boano	&	gloss	&		\\	\midrule
%\endfirsthead
\tbx{10-32}	&	PMP	&	Boano	&	gloss		\\	\midrule
%\endhead
	&	*\tbr{R}umbia	&	\tbr{l}ipia	&	sago			\\
	&	*saku\tbr{R}u	&	saʔa\tbr{l}u	&	reef (irregular *k > ʔ)	\\
	&	*ba\tbr{R}ah	&	au ha\tbr{l}a	&	ember			\\
	&	*na\tbr{R}ah/*na\tbr{r}a	&	na\tbr{l}a	& tree, \it{Ptereocarpus indicus}	\\
	&	*\tbr{R}ibut	&	\tbr{l}ihut-e	&	storm			\\
\lspbottomrule
\end{tabular}\mbox{}\\}

The reflexes of PMP *l are quite complex,
splitting *l > \it{l, n}.
There is a retention of initial *l = \it{l},
as shown in example \qf{10-33}.
\\

\noindent{\wg\st{5pt}\begin{tabularx}
{\linewidth}{>{\hspace{-\tabcolsep}}l<{\hspace{-\tabcolsep}}
>{\hspace{-\tabcolsep}\enspace}lll
>{\raggedright\arraybackslash}X}
\lsptoprule
%\tbx{10-33}	&	PMP/PCM	&	Boano	&	gloss	\\	\midrule
%\endfirsthead
\tbx{10-33}	&	PMP/PCM	&*gloss&	Boano	&	gloss	\\	\midrule
%\endhead
	&	*\tbr{l}uaq		&spit out, spew				&	\tbr{l}ua				&	vomit	\\
	&	*\tbr{l}umut	&moss									&	\tbr{l}umut-e		&	moss	\\
	&	*\tbr{l}umpuk	&cluster, group				&	paka-\tbr{l}upu	&	gather together, collect	\\
	&	*\tbr{l}asəR	&scrotum and testicles&	\tbr{l}ase			&	scrotum and testicles	\\
	&	*\tbr{l}awan	&oppose								&	\tbr{l}aa				&	opposite, other side	\\
	&	*\tbr{l}imas	&bail out							&	\tbr{l}ima-te		&	bailer (n)	\\
	&	**\tbr{l}abu	&ashes, dust					&	\tbr{l}ahu-ʔi		&\hp{\,}	\\ \hhline{~}
	& < *qabu 			&	\hp{\,}							&	\hp{\,}					&\hp{\,}	\\ \hhline{~}
	&\hp{\,}				&\hp{\,}							&\hp{\,}					&\hp{\,}	\\ \hhline{~}
	&\hp{\,}				&\hp{\,}							&\hp{\,}					&\multirow{-4}{118pt}{ashes, dust;
																				sporadic *l onset in several languages,
																				cf. PMP *dabuk ‘ashes’ \citep[778]{Wolff2010}} \\
\lspbottomrule
\end{tabularx}\mbox{}\\}

Initial *l > \it{n} is through regressive assimilation where
/n/ is the next consonant,
as shown in example \qf{10-34}.

\noindent{\wg\st{6pt}\tblr{>{\hspace{-\tabcolsep}}l<{\hspace{-\tabcolsep}}
>{\hspace{-\tabcolsep}\enspace}
lll}
\lsptoprule
%\tbx{10-34}	&	PMP	&	Boano	&	gloss	\\	\midrule
%\endfirsthead
\tbx{10-34}	&	PMP	&	Boano	&	gloss	\\	\midrule
%\endhead
	&	*ma-\tbr{l}inaw	&	\tbr{n}ina	&	calm (water) 	\\
	&	*\tbr{l}iaŋ	&	\tbr{n}ian-e	&	cave	\\
	&	**qa\tbr{l}ipan	&	ni{\tl}\tbr{n}ian-e	&	centipede	\\
	&	*\tbr{l}aŋit	&	\tbr{n}anit-e	&	sky	\\
	&	*\tbr{l}ayaR	&	\tbr{n}aan	&	sail (n.)	\\
\lspbottomrule
\end{tabular}\mbox{}\\}

\newpage
Word medially in proximity to high vowels *l > \it{n},
as shown in \qf{10-35}.
Between non-high vowels *l = \it{l} as shown in \qf{10-36}.
The different reflexes of *l adjacent to high vowels are reminiscent
of the patterns discussed in \srf{10-01}.
\vspace{2pt}

\noindent{\wg\st{6pt}\begin{tabularx}
{\linewidth}{>{\hspace{-\tabcolsep}}l<{\hspace{-\tabcolsep}}
>{\hspace{-\tabcolsep}\enspace}ll
>{\raggedright\arraybackslash}X}
\lsptoprule
%\tbx{10-35}	&	PMP/PCM	&	Boano	&	gloss	\\	\midrule
%\endfirsthead
	&	PMP/PCM	&	Boano	&	gloss	\\	\midrule
%\endhead
\tbx{10-35}	&	*pi\tbr{l}iq	&	para-i\tbr{n}i-ʔe	&	choose	\\
	&	*bu\tbr{l}an	&	hu\tbr{n}an-e	&	moon	\\
	&	*qu\tbr{l}u	&	u\tbr{n}u	&	head	\\
	&	*bu\tbr{l}u	&	hu\tbr{n}u	&	body hair	\\
	&	*tə\tbr{l}u	&	te\tbr{n}u	&	three	\\
	&	**batə\tbr{l}u	&	hate\tbr{n}u	&	thick (cf. \ili{Hitu} \it{hatelu}, \ili{Larike} \it{hutido}, \ili{Geser} \it{batolu} among others; metathesis from PMP *təbəl)	\\
	&	*qahə\tbr{l}u	&	a\tbr{n}u	&	pestle	\\
	&	*qu\tbr{l}u	&	u\tbr{n}u	&	head	\\
	&	*bu\tbr{l}an	&	hu\tbr{n}an-e	&	moon	\\
	&	*qa\tbr{l}ima	&	\tbr{n}ima	&	hand	\\ \midrule
\tbx{10-36}	&	*ka\tbr{l}aw	&	a\tbr{l}a-wele	&	hornbill	\\
	&	*da\tbr{l}əm	&	-la\tbr{l}e	&	in, inside (*l = \it{l} could also be due to progressive assimilation)	\\
	&	*za\tbr{l}an	&	la\tbr{l}an-e	&	path (*l = \it{l} could also be due to progressive assimilation)	\\
\lspbottomrule
\end{tabularx}\mbox{}\\}

The word \it{nima} ‘hand’ appears to be an exception with initial
\it{n}, and contrasts with the \it{lima-te} ‘bailer’ in example \qf{10-33}.
But note that historically,
*l in *qalima > \it{nima} ‘hand’ is a medial consonant,
whereas in *limas > \it{lima-te} it is an initial consonant.
If this analysis is correct,
then the trace of antepenultimate PMP *qa- is significant,
and is one of several examples showing that PMP *qa- cannot be
said to be completely lost in the region (contra \citealp[263f]{Blust1993}).

Three exceptions have been found in which *l = \it{l} is adjacent
to a high vowel.
The word for ‘garfish’ is in a semantic domain that could easily be a loan.
The word for ‘wart’ is historically trisyllabic
and has *l in close proximity to *R.

\noindent{\wg\st{6pt}\tblr{>{\hspace{-\tabcolsep}}l<{\hspace{-\tabcolsep}}
>{\hspace{-\tabcolsep}\enspace}
lll}
\lsptoprule
%\tbx{10-37}	&	PMP	&	Boano	&	gloss	\\	\midrule
%\endfirsthead
\tbx{10-37}	&	PMP	&	Boano	&	gloss	\\	\midrule
%\endhead
	&	*sə\tbr{l}u	&	se\tbr{l}u	&	longtom, garfish	\\
	&	*ti\tbr{l}u/*tu\tbr{l}i	&	ti\tbr{l}u	&	PMP ‘earwax, deaf’; Boano ‘earwax’	\\
	&	*butə\tbr{l}iR	&	huti\tbr{l}-e	&	wart	\\
\lspbottomrule
\end{tabular}\mbox{}\\}

\citet[69]{Collins1983} posited a close relationship between
Boano, \ili{Larike}, and \ili{Wakasihu} and placed them in a single branch (\tsc{East
Hoamoal}) of \tsc{West \ili{Piru} Bay}.
We note that the conditioning environment for PMP *l > \it{n}
(adjacent to a high vowel) is the same as that for \ili{Larike} *l > \it{d}.
However, this is also the conditioning environment for **l >
\it{r} in \ili{Amahai}, \tsc{Kamarian}, \ili{Kaibobo}, and \ili{Haruku}, which \citet[100]{Collins1983}
places within \tsc{East \ili{Piru} Bay}.

Like *l, PMP *j > \it{n} adjacent to high vowels,
but *j > \it{l} adjacent to non-high vowels.
There is also one case of *j > \it{n} between non-high vowels
which is likely a case of progressive assimilation to initial *n.
These examples are shown in \qf{10-38}.
Word finally *j > {\0},
though there is only one example in our data: *waRəj > \it{ane-te}
‘vine, rope’.

\noindent{\wg\st{5.5pt}\tblr{>{\hspace{-\tabcolsep}}l<{\hspace{-\tabcolsep}}
>{\hspace{-\tabcolsep}\enspace}
llll}
\lsptoprule
%\tbx{10-38}	&		&	PMP	&	Boano	&	gloss	\\	\midrule
%\endfirsthead
\tbx{10-38}	&		&	PMP	&	Boano	&	gloss	\\	\midrule
%\endhead
	&	*j > \it{n} /V\tsc{\tsc{[+high]}}	&	*qapə\tbr{j}u	&	ne\tbr{n}u	&	gall, bile (unexplained initial \it{{\#}n})	\\
	&		&	*ŋi\tbr{j}uŋ	&	ni\tbr{n}u	&	nose	\\
	&		&	*pi\tbr{j}a	&	i\tbr{n}a	&	how many	\\
	&		&	*hua\tbr{j}i	&	a\tbr{n}i	&	younger sibling	\\
	&	*j > \it{l} /V\tsc{\tsc{[-high]}}	&	*pa\tbr{j}ay	&	a\tbr{l}a	&	rice	\\
	&		&	*ma\tbr{j}a	&	ka-ma\tbr{l}a	&	dry up	\\
	&	*j > \it{n} /\it{n}V{\gap}	&	*ŋa\tbr{j}an	&	na\tbr{n}a	&	name	\\
\lspbottomrule
\end{tabular}\mbox{}\\}

We see the unmarked reflex of PMP *l/*j in Boano as \it{l},
with *l/*j > \it{n} as a conditioned reflex
(/V\tsc{\tsc{[+high]}} and/or progressive assimilation).
Recall, however, that *R > \it{n} in all environments,
not just adjacent to high vowels or as cases of assimilation,
as shown in \qf{10-29}.

In our interpretation, the Boano data show *R > \it{n},
and *l/*j > \it{l},
with only the conditioned reflex of *l/*j > \it{n} merging with *R > \it{n}.
To put it another way,
the base correspondences of PMP *R
and *l/*j are different,
and the environments of the conditioned reflexes of PMP *R
and *l/*j are different from each other,
so we cannot support a merger of PMP *l
and *R for the language ancestral to Boano.
The reflexes of *d/*z/*R/*l/*j in Boano under our account are
summarised in \qf{10-BoanoChange}, leaving aside the sporadic exceptions:

\noindent{\st{6pt}\tblr{>{\hspace{-\tabcolsep}}l<{\hspace{-\tabcolsep}}
>{\hspace{-\tabcolsep}\enspace}
ll@{ }l@{ }ll}
\tbx{10-BoanoChange}	& a.	& *d/*z	&>&	\it{l}	&	(all positions)\\
											& b.	& *R		&>&	\it{l}	&	/{\#}lV{\gap}\\
											& c.	& *R		&>&	\it{n}	&	elsewhere\\
											& d.	& *l/*j	&>&	\it{l}	&	/V\tsc{[-high]}{\gap}V\tsc{[-high]}\\
											& e.	& *l/*j	&>&	\it{n}	&	/V\tsc{[+high]}\\
\end{tabular}\mbox{}\\}

%\begin{tabular}{ll@{ }l@{ }ll}
	%1. & *d/*z	&>&	\it{l}	&	(all positions)\\
	%2. & *R			&>&	\it{l}	&	/{\#}lV{\gap}\\
	%3. & *R			&>&	\it{n}	&	elsewhere\\
	%4. & *l/*j	&>&	\it{l}	&	/V\tsc{[-high]}{\gap}V\tsc{[-high]}\\
	%5. & *l/*j	&>&	\it{n}	&	/V\tsc{[+high]}\\
%\end{tabular}\mbox{}\\

The implication of all this Boano data is two-fold.
Firstly, if Boano truly belongs within \tsc{\ili{Piru} Bay},
then Boano split off before the merger of PMP *d/*z/*R/*l/*j
that occurs in \ili{Larike}
and \ili{Wakasihu} (see \srf{10-01}).
Secondly, the merger of PMP *R/*l/*j > **l cannot be assigned
to the level of Proto-\ili{Ambon}-Seram.
With regard to these segments,
a simpler treatment might be to have Boano be its own branch
within \tsc{\ili{Ambon}-Seram} rather than linked to \ili{Larike}
and \ili{Wakasihu}, with *R/*l/*j > **l occurring in other \tsc{\ili{Ambon}-Seram} languages.
If further work shows a close relationship between Boano,
\ili{Larike}, and \ili{Wakasihu}, then we would propose that the branch
containing these languages (\citeauthor{Collins1983}' \tsc{East Hoamoal}) be a primary
branch of \tsc{\ili{Ambon}-Seram}.

If Boano is a primary branch,
the change *l/*j > **l can be assigned to the level of \tsc{\ili{Ambon}-Seram}
with Boano then undergoing **l > \it{n} /V\tsc{[+high]} followed
by a merger of *d/*z
and other instances of **l to \it{l} in other positions.
The changes affecting *R in Boano can be left unordered with
respect to these changes.

\citet[95--97]{Collins1983} sees that the Boano data
and the \ili{Asilulu} data together “indicate retention of a reflex
of *ɫ [tilde l] which in all positions is distinct from *l.
These two retentions must be acknowledged to have occurred in
the proto-language of these two languages.
Hence, our understanding of the inventory of the proto-sounds
in Proto-\ili{Central} Maluku must be revised.
Indeed, the \ili{Asilulu} evidence presented here impels us to make
another revision in our notion of the sounds of non-Formosan
languages, Blust’s Proto-Malayo-Polynesian.” Updating the words given by
Collins in support of *ɫ,
we find that they all currently have forms (or doublets) with
*n, so there is no need to appeal to a distinction between *ɫ
and *l, nor to posit an additional sound in the proto inventory.\footnote{
		The words for which \citet[96--97]{Collins1983} posits *ɫ include:
		PMP *naŋuy ‘swim’, *nuka ‘wound, sore’, *niŋal ‘echo’, *qaninu ‘shadow’,
		and *anak ‘child’.} 
